%\section{Additional material on \piI }
%\label{a:pi}







              





% \subsection{Equivalences and preorders }


% \finish{
% i could talk about barbed congruence in place of bisimilarity (as in the paper with
% Adrien) because the encodings are image-finite upto $\wbPI$
% I should also say that all relations are defined on configurations with the same $\phi$ 
% (should these be still called processes?)
% } 

% \begin{definition}
% \label{d:traces}
% A \emph{trace}  is  a finite  sequence of visible actions; we use 
%  $\epsilon $ for the
%  empty   sequence.  Thus the grammar for traces is
% \[\sigma : = \epsilon \midd \mu; \sigma  \]
% If $\sigma $ is a trace, then $P \ArrPI \sigma $ is defined inductively thus: 
% \begin{itemize}
% \item $P \ArrPI \epsilon P$.
% \item $P \ArrPI {\mu;  \sigma }$ if there is $P'$ with $P \ArrPI\mu P'$ and 
% $P' \ArrPI\sigma $.
% \end{itemize} 
% Processes $P$ and $Q$ are  \emph{trace equivalent} 
% %, written   $P \may Q$ if $\trace P  = \trace Q$.
% if 
% $P \ArrPI\sigma $  implies $Q \ArrPI\sigma $,  and conversely, 
% for all traces $\sigma$.
% \end{definition} 
% We write
% \begin{itemize}
% \item  $\equiv$ for structural congruence;
% \item $\wbPI$ for weak bisimilarity; 
% \end{itemize} 
% \finish{to be completed, with the definitions of trace eq, weak bisimulation, structural
%   congruence (maybe struct congr is not needed, though) } 


 

\section{Auxiliary material for Section~\ref{subsec:seq-ts-pi}}
%\section{Behavioural relations in the ordinary LTS and  in the \opLTS for \piI } 
\label{a:behav}

\DSOCT{ moved the first half of this appendix onto the main text, and consequently changed
title of the appendix}



% A behavioural relation on the ordinary LTS `almost always' implies the corresponding
% relation on the \opLTS.  The only exceptions are produced by processes that are related
% although they exhibit different divergence behaviours (by divergence we mean the
% possibility of performing an infinite number of $\tau$ actions) and at the same time may
% perform input actions. For instance, using $\tau^\omega$ to indicate a purely divergent
% process  (i.e., a process whose only transition is $\tau^\omega \arr\tau \tau^\omega$), 
%  the processes $a.\nil | \tau^\omega $ and $a.\nil$ are  bisimilar (and trace) equivalent
%  in the ordinary LTS, but they are not so in the \opLTS, since the divergence in the former
%  process prevents observation of the input at $a$. 


% Indeed, a `bisimilarity respecting divergence' would imply bisimilarity on the \opLTS. 
% Such  `bisimilarity respecting divergence', written $\wbPIDIV$, is defined as $\wbPI$ with
% the additional requirement that $P \wbPI Q$ implies 
% \begin{itemize}
% \item
% $P \Uparrow $ iff  $Q \Uparrow $ \hfill $(*)$
% \end{itemize} 
% where the predicate $ \Uparrow $ holds for an  agent $R$  if it
% has a divergence (an infinite path 
% consisting of only $\tau$ actions). 

% \begin{lemma}
% \label{l:bisDIV_op}
% Relation $\wbPIDIV$ implies $\wbOPI$.
% \end{lemma} 
% The same result applies to  other relations; e.g., the `expansion  respecting divergence'
% (defined by adding the clause $(*)$ above to those of the expansion
% relation $\expa$)
%  on the ordinary LTS implies expansion  on the \opLTS.

% The results of operational correspondence between \piI and \AOGS rely on
% Lemmas~\ref{l:opt_sound}-\ref{l:opt_con}. In the proofs of all these lemmas
% \cite{Durier20}, the uses of  expansion ($\expa$) come from application of some
% simple algebraic laws that respect  divergence (i.e., the laws do not add or remove any
% divergence). Indeed the laws either are valid for strong bisimilarity, or they are laws
% such as 
% \[ 
% {\res{a} ( \inp a {\tilx}.P | \bout a \tilx . Q)} \;  \contr \;
% {\res{a} (P |  Q)  }
%  \]
% As a consequence, all occurrences of $\expa$ in Lemmas~\ref{l:opt_sound}-\ref{l:opt_con}
% can be replaced by the expansion relation defined on the \opLTS ($\expaOP$). 
% This in turn implies that the same could be done for 
%  all occurrences of $\expa$ in
% Theorems~\ref{t:opcorrCON_w}-\ref{t:opcorrCON3_w}   (that rely on the lemmas).




On
 the encoding of $\lambda$-terms, 
the ordinary
bisimilarity $\wbPI$ implies the `bisimilarity respecting divergence' $\wbPIDIV$. 
%
%
% First we recall 
% that the equivalence induced on call-by-value  $\lambda$-terms by their encoding into \piI 
% coincides with 
% Lassen's 
% \emph{eager normal-form (enf) bisimilarity} \cite{lassentrees,lassenfa} (also recalled in
% Definition~\ref{d:enf}).
% \begin{theorem}[\cite{DurierHS18}]
% \label{t:adrien}
% $\enca M \wbPI \enca N$ iff $M$ and $N$ are enf-bisimilar. 
% \end{theorem} 
% On such terms, 
% the ordinary
% bisimilarity $\wbPI$ implies the `bisimilarity respecting divergence' $\wbPIDIV$. 
This
because, intuitively, a term  $\encoIa {M} p$ is divergent iff $M$ is so in the call-by-value
$\lambda$-calculus.    Formally, the result is proved by rephrasing the result about the 
correspondence between $\wbPI$ and Lassen's trees (Theorem~\ref{t:adrien}) in terms of 
$\wbPIDIV$ in place of $\wbPI$. Again, this property boils down to the fact that the uses of
$\expa$ in Lemmas~\ref{l:opt_sound}-\ref{l:opt_con} (on which also the characterisation in
terms of Lassen's trees is built) respect divergence.  
Thus, using also Lemma~\ref{l:bisDIV_op}, we derive the following result. 

\begin{theorem}
\label{t:wbPI-wbOPI}
For $M,N \in \Lambda $, we have:
$\encoIa {M} p \wbPI \encoIa {N} p$ iff 
$\encoIa {M} p \wbOPI \encoIa {N} p$. 
\end{theorem} 

% (Corollary~\ref{c:} shows that also the converse holds; in general however $\wbOPI$ does
% not imply $\wbPI$ or even $\wbPIDIV$.)



\ds{The 7.1 should be rewritten for \opLTS -- at least put in in the appendix} 


\ds{below i talk about ``singleton configurations'', not specifying Alternating or COGS,
  for instance, because the grammar of singleton conf is the same and also then the
  encoding. Not sure if we want to say something}


\noindent
Below we report an alternative proof of the theorem above, however
always relying  on the variant of lemmas, such as  
 Lemma~\ref{l:opt_sound} and Lemmas~\ref{l:opt_stuck}-\ref{l:opt_con}, with 
$\expaOP $ in place of $\expa$. 

\begin{lemma}
\label{l:shapesT}
If $\FF$ is a singleton configuration, then  $\encoConI{\FF}$ is of three possible forms: 
\begin{enumerate}
\item 
$\encoIa {M} p$, for some $M,p$;
\item
$\inp qx .\encoIa {E[x]} p $, for some $E,x,q,p$;
\item 
$\encoIVa {V} y$, for some $V,y$ (that is, either  $!\inp y {x,q}.\encoIa M q$ if $V  =
\lambda x. M$, or $ \fwd y x$ if $V = x$).
\end{enumerate}  
\end{lemma} 


\begin{lemma}
\label{l:deco}
Suppose  $a$ not free in $P_1 $ and $Q_1$. 
\begin{itemize}
\item
Let $P = P_1 | ! \inp a {\tilb}.P_2$, 
and $Q = Q_1 | ! \inp a {\tilb}.Q_2$, 
and  $ P \wbPI  Q$.
Then also 
  $ {P_1}  \wbPI  {Q_1}$.


\item 
Similarly, let $P = P_1 |  \inp a {\tilb}.P_2$, 
and $Q = Q_1 |  \inp a {\tilb}.Q_2$, 
and  $ {P} \wbPI  {Q}$.
Then also 
  $ {P_1}  \wbPI  {Q_1}$.
\end{itemize}
% Suppose $P=  P_1|P_2$ and 
% $Q = Q_1|Q_2$, with
% either 
%  ($P_1 $ of the form $\inp a \tilx . P'$ and 
%  $Q_1 $ of the form $\inp a \tilx . Q'$)
% or  
%  ($P_1 $ of the form $!\inp a \tilx . P'$ and 
%  $Q_1 $ of the form $!\inp a \tilx . Q'$)
% for some  $a$ not free in $P_2 $ and $Q_2$. 
% Then $P \wbPI Q $ implies $P_2 \wbPI Q_2$.
\end{lemma} 

\begin{theorem}
\label{t:opBIS_BIS}
If $\FF$ and $\GG$ are singleton configurations, then
$\encoConI  \FF \wbPI \encoConI  \GG $  implies 
 $\encoConI  \FF \wbOPI \encoConI  \GG $.
\end{theorem} 

\begin{proof}
Let $\R$ be the relations with all pairs $(P,Q)$ where $P$ and $Q$ are of the
form 
\[
\begin{array}{rcl}
 P &=&   %\encoConI{\FF_1} | \ldots |  \encoConI{\FF_n} 
P_1  | \ldots | P_n
\\
 Q &=& %\encoConI{\GG_1} | \ldots |  \encoConI{\GG_n} 
Q_1  | \ldots | Q_n
\end{array}
 \] 
for some $n$, where all $P_i,Q_i$ are 
encodings of singleton configurations and  for all $i, $ we have $P_i \wbPI Q_i$.
Moreover:
\begin{itemize}
\item
 there is at most one $i$ for which  $P_i,Q_i$ are encodings of  active
configurations;
\item  in any $P_j$ (resp.\ $Q_j$) that is the encoding of a passive
configuration, the name of the initial input does not appear free in any other component
$P_{j'}$ (resp.\ $Q_{j'}$) for $j'\neq j$. 
\end{itemize} 
\noindent
A consequence of the two conditions above is that two distinct components $P_i$ and $P_j$
cannot interact. That is, if $P \arr \mu P'$,  then there is $i$ such
that $P_i \arr \mu P'_i$ and $P' = P''_1| \ldots | P''_n$ where $P''_j = P_j $ if $j\neq i$
and $P''_i = P'_i $. And similarly if   $P \arrN \mu P'$.



We  show that  $\R$ is a  $\wbOPI$-bisimulation  up-to $\expaOP$ (the bisimulation
up-to expansion is sound for bisimilarity in any LTS). 
In the proof below,  Lemmas~\ref{l:opt_op}  and 
\ref{l:opt_con}
actually refer to the versions of the lemmas mentioned above, with $\expaOP$
in place of $\expa$.  


Suppose $P \arrN \mu P'$.  The action orginates from some $P_i$ alone, say $P_i \arrN\mu P_i'$.
If $\mu$ is a $\tau $ action, then also $P_i \arr\mu P_i'$; 
and since $P_i \wbPI Q_i$, we have $Q_i \Longrightarrow  Q_i'$ with $P_i' \wbPI Q_i$.  
Moreover, also $Q_i \RaN  Q_i'$ --- and the corresponding transition from $Q$.


Using Lemma~\ref{l:opt_op} we infer that, up-to expansion, the
derivative processes $P_i'$ and $Q'_i$ can be rewritten into processes
that fit the definition of $\R$. 

For the case when $\mu$ is an output, one reasons similarly, possibly
also using Lemma~\ref{l:deco} when $\mu $ is of the form $\bout x {z,q}$.  

Finally, if $P \arrN \mu P'$ and $\mu$ is an input, then $P $ is input
reactive. From the conditions in the definition of $\R$ we infer that
$Q$ is input reactive too.  Then we can conclude, reasoning as in the
previous cases, but this time  using Lemma~\ref{l:opt_con} with
$\gamma   $ a singleton.  
\end{proof}  
\noindent
Theorem~\ref{t:wbPI-wbOPI}  is then a corollary of
Theorem~\ref{t:opBIS_BIS}.  


% However
% in the processes considered in the paper (encoding of %$\lambda$-terms and of 
% OGS
% configurations) the implication does hold: a behavioural relation on the ordinary LTS
% implies the corresponding relation on the \opLTS. 
% %We omit  the details because the results are actually not needed in the paper. 
% We report the result  below.
% The intuitive explanation is the following. 
% In  the two  LTSs, outputs and internal actions are the same.  The only constraint concerns
% inputs, for the \opLTS, where inputs are visible only on  input reactive processes.
% The only input reactive processes in the paper are those encoding passive configurations, 
% and  such processes are only related to  themselves (i.e., to encodings of passive
% configurations).
% %  moreover,
% % such processes  arise from (encodings of)  active configurations and of  running $\lambda$-terms
% % precisely after an  
% % output action (see the results in Appendix~\ref{a:oc} such as Lemmas~\ref{l:opt_op} and
% % \ref{l:opt_con}). 

% In fact, on the processes of the paper, all behavioural relations considered, both in the
% ordinary LTS and in the \opLTS, coincide. This can be deduced from
% Corollaries~\ref{c:fa_traces_bis}, \ref{c:ogsp_pi_traces_bis}, and
% \ref{c:OGS_COGS}.  


